\section{Proofs for Finite Calibration}
\label{sec:app-calibration}

\subsection{Regularity of the Target Covariance Transform}

For \cref{prop:regularity}, if
$g-\E g=\sum_{k\ge1}(a_k/k!)H_k$, then
$C_g(r)=\sum_{k\ge1}(a_k^2/k!)r^k$. For
$g\in W^{1,2}(\phi)$, termwise differentiation and the Gaussian Sobolev
identity give
\begin{equation}
L_g:=\sup_{|r|\le1}|C_g'(r)|
\le \sum_{k\ge1}\frac{k a_k^2}{k!}
=\E\{g'(U)^2\}<\infty.
\label{eq:app-smooth-L}
\end{equation}
The non-negative coefficients show that equality in the supremum is attained
at $r=1$, proving the Lipschitz claim in part (i).

For $g_c(z)=\Ind\{z>c\}$, Plackett's identity gives, for $-1<r<1$,
\begin{equation*}
C_{g_c}'(r)=
\frac{\exp\{-c^2/(1+r)\}}{2\pi\sqrt{1-r^2}}.
\end{equation*}
Together with the sharp arcsine inequality
$|\arcsin b-\arcsin a|\le(\pi/\sqrt2)|b-a|^{1/2}$, this yields the global
bound
\begin{equation}
|C_{g_c}(r_2)-C_{g_c}(r_1)|
\le \frac{1}{2\sqrt2}|r_2-r_1|^{1/2}.
\label{eq:app-threshold-holder}
\end{equation}
At $c=0$ and $(r_1,r_2)=(-1,1)$ the constant is attained. Expanding the same
derivative as $r\uparrow1$ and integrating over $[1-\delta,1]$ gives
\eqref{eq:boundary-asymptotic}, so no larger global H\"older exponent is
possible. Finally, on $[-r_{\max},r_{\max}]$ the derivative is bounded by
\begin{equation}
L_{g_c}(r_{\max})=
\frac{\exp\{-c^2/(1+r_{\max})\}}
{2\pi\sqrt{1-r_{\max}^2}},
\label{eq:app-threshold-local-L}
\end{equation}
which proves part (iii).

\subsection{Covariance Repair and Resolvent Bounds}

For covariance repair, put $\bar K=\widehat\Sigma-\widehat R_0$ and
$e_0=\opnorm{\widehat\Sigma-(K+R_0)}+\opnorm{\widehat R_0-R_0}$.
If $0<\tau\le\lambda_{\min}(K)/2$ and
$e_0\le\min\{1/2,\lambda_{\min}(K)/2\}$, Weyl's inequality gives
$\bar K\succeq\tau I_p$, so flooring is inactive. For
$H=\Diag\{\diag(\bar K)\}^{-1/2}$, the unit diagonal of $K$ gives
$\opnorm{H-I_p}\le2e_0$ and $\opnorm H\le2$. Expanding
$H\bar K H-K=(H-I_p)\bar K H+(\bar K-K)H+K(H-I_p)$ yields
\begin{equation}
\opnorm{\cR_\tau(\bar K)-K}
\le(6\opnorm K+4)e_0
\le6(1+\opnorm K)e_0.
\label{eq:app-repair-bound}
\end{equation}
Thus diagonal rescaling preserves the raw rate at an interior $K$.

For the uniform resolvent bound, let $E=\widehat K-K$, $e=\opnorm E$, and
$\kappa=\max\{\opnorm K,\opnorm{\widehat K}\}$. Under
\eqref{eq:family-conditioning}, if $ae\le\lambda/2$, then for every non-empty
$S\in\Pi_{\mathcal B}$ the inverse and resolvent bounds are
\begin{equation*}
\begin{aligned}
\opnorm{(A_S\widehat K A_S^\top+R_S)^{-1}}&\le2/\lambda,\\
\opnorm{(A_S\widehat K A_S^\top+R_S)^{-1}
-(A_SKA_S^\top+R_S)^{-1}}&\le2ae/\lambda^2.
\end{aligned}
\end{equation*}
Expanding the two outer factors and the inverse in $Q_S$ then yields
\begin{equation}
\opnorm{Q_S(\widehat K)-Q_S(K)}\le L_Qe,
\qquad
L_Q=\frac{4\kappa a}{\lambda}
    +\frac{2\kappa^2a^2}{\lambda^2},
\label{eq:app-resolvent-bound}
\end{equation}
uniformly over the feasible family. Thus the constant worsens only through the
latent covariance scale, measurement amplification $a$, and protocol
conditioning $\lambda$.

\subsection{Value Error, Calibration Rates and Regret}

For \cref{thm:uniform-error}, the empty protocol has value zero under both
matrices. For non-empty $S$, non-negative weights summing to one, the
target modulus and \eqref{eq:app-resolvent-bound} imply
\begin{equation*}
|F_g(S;\widehat K)-F_g(S;K)|\le L(L_Qe)^\beta,
\qquad
|V_g(\widehat K)-V_g(K)|\le Le^\beta.
\end{equation*}
For the local threshold case, the second bound uses that $K$ and
$\widehat K$ have the same unit diagonal, so only off-diagonal arguments vary.
Since $0\le F_g(S;K)\le V_g(K)$, the ratio identity gives the more explicit
uniform bound
\begin{equation}
\sup_{S\in\Pi_{\mathcal B}}
|\widehat\cI_g(S)-\cI_g(S)|
\le
\frac{L(L_Q^\beta+1)e^\beta}{V_g(K)-Le^\beta}.
\label{eq:app-explicit-uniform}
\end{equation}
For sufficiently small $e$, the denominator is at least $v_0/2$ and
$\kappa$ is bounded by a fixed neighbourhood of $K$. Hence
\eqref{eq:app-explicit-uniform} is at most
$2L(L_Q^\beta+1)e^\beta/v_0$, which is \eqref{eq:uniform-error} with exactly
the dependencies stated there. For a threshold target satisfying
\eqref{eq:r0}, continuity of all finitely or uniformly conditioned resolvents
keeps the varying correlations in a common compact subinterval of $(-1,1)$;
the local Lipschitz constant \eqref{eq:app-threshold-local-L} then applies.

For \cref{thm:fixed-model}, fixed-dimensional Gaussian covariance
concentration gives
\begin{equation*}
\opnorm{\widehat\Sigma-(K+R_0)}=O_p(m^{-1/2}).
\end{equation*}
The assumed rate for $\widehat R_0$ therefore makes
$e_0=O_p(m^{-1/2})$. Because $K\succ0$ and
$\tau_m\downarrow0$, the conditions of \eqref{eq:app-repair-bound} hold with
probability tending to one, giving
$\opnorm{\widehat K-K}=O_p(m^{-1/2})$.

Substitution into \eqref{eq:app-explicit-uniform} gives root-$m$ value error
when $\beta=1$ and $m^{-1/4}$ under the global threshold envelope
$\beta=1/2$. At a fixed threshold model satisfying \eqref{eq:r0}, the local
Lipschitz modulus has $\beta=1$, restoring the root-$m$ rate. The family is
finite, so the conditioning and correlation neighbourhoods can be chosen
uniformly over all $S\in\Pi_{\mathcal B}$, which gives the stated rate.

For \cref{cor:regret}, insert the empirical value on both sides of the
comparison between $S^*$ and $\widetilde S$:
\begin{align*}
\cI_g(S^*)-\cI_g(\widetilde S)
&\le
\{\cI_g(S^*)-\widehat\cI_g(S^*)\}
+\{\widehat\cI_g(S^*)-\widehat\cI_g(\widetilde S)\}
+\{\widehat\cI_g(\widetilde S)-\cI_g(\widetilde S)\}\\
&\le2\varepsilon_m+\eta_m,
\end{align*}
which proves \eqref{eq:regret-opt}. Applying the same argument within
$\Pi^{(\ell)}$ gives
$\cI_g(\widehat S_\ell)\ge
\max_{S\in\Pi^{(\ell)}}\cI_g(S)-2\varepsilon_\ell$; subtracting this lower
bound from the global optimum gives \eqref{eq:resolution-bound}.

\subsection{Best-Linear Calibration}
\label{sec:app-best-linear}

\theoremstyle{plain}
\newtheorem*{bestlinearproposition}{Proposition}
\begin{bestlinearproposition}[Uniform root-$m$ error]
For non-empty $S$, write $\Sigma_S=\Var(Y_S)$ and
$c_S=\Cov(Y_S,\Theta)$, and let $v=\Var(\Theta)$. Suppose, uniformly over
$S\in\Pi_{\mathcal B}$, that $v\ge v_0>0$,
$\lambda_{\min}(\Sigma_S)\ge\lambda>0$, $\|c_S\|_2\le C_c<\infty$, and
\begin{equation*}
\delta_m:=\sup_{S\ne\varnothing}
\{\opnorm{\widehat\Sigma_S-\Sigma_S}+\|\widehat c_S-c_S\|_2\}
+|\widehat v-v|=O_p(m^{-1/2}).
\end{equation*}
Set both values of the empty protocol to zero; otherwise use
$\cI_{\mathrm L}(S)=c_S^\top\Sigma_S^{-1}c_S/v$ and its plug-in version.
Then
$\sup_{S\in\Pi_{\mathcal B}}|\widehat\cI_{\mathrm L}(S)-\cI_{\mathrm L}(S)|
=O_p(m^{-1/2})$.

In the measurement model, where $\Sigma_S=A_SKA_S^\top+R_S$ and
$c_S=A_Sc$ for $c=\Cov(Z,\Theta)$, it suffices that
$\sup_S\opnorm{A_S}<\infty$ and
\begin{equation*}
\opnorm{\widehat K-K}+\|\widehat c-c\|_2+|\widehat v-v|
+\sup_S\opnorm{\widehat R_S-R_S}=O_p(m^{-1/2}),
\end{equation*}
using $\widehat\Sigma_S=A_S\widehat K A_S^\top+\widehat R_S$ and
$\widehat c_S=A_S\widehat c$.
\end{bestlinearproposition}

\begin{proof}
On the event $\delta_m\le\min(\lambda/2,v_0/2)$, the resolvent identity
and $q_S=c_S^\top\Sigma_S^{-1}c_S$ give, uniformly over non-empty $S$,
\begin{equation*}
\begin{aligned}
\opnorm{\widehat\Sigma_S^{-1}}&\le\frac{2}{\lambda},
&\qquad
\opnorm{\widehat\Sigma_S^{-1}-\Sigma_S^{-1}}
&\le\frac{2\delta_m}{\lambda^2},\\
\sup_S|\widehat q_S-q_S|&\le
\frac{2\delta_m}{\lambda}(2C_c+\delta_m)
+\frac{2C_c^2\delta_m}{\lambda^2}
=O_p(m^{-1/2}).
\end{aligned}
\end{equation*}
Because $q_S\le C_c^2/\lambda$ and $\widehat v\ge v_0/2$, division by
$\widehat v$ and $v$ proves the claim. The primitive rate implies the assumed
one through
$\opnorm{\widehat\Sigma_S-\Sigma_S}
\le\opnorm{A_S}^2\opnorm{\widehat K-K}
+\opnorm{\widehat R_S-R_S}$ and
$\|\widehat c_S-c_S\|_2\le\opnorm{A_S}\|\widehat c-c\|_2$.
\end{proof}
